% ============================================================
%  CATEGORICAL SEARCH FOR METABOLOMICS ANNOTATION:
%  GENERATIVE PHASE-SPACE ADDRESSING WITH COHERENCE-GROUNDED
%  VERIFICATION FOR UNKNOWN COMPOUND IDENTIFICATION
% ============================================================

\documentclass[twocolumn,10pt,a4paper]{article}

\usepackage[utf8]{inputenc}
\usepackage[T1]{fontenc}
\usepackage{lmodern}
\usepackage[a4paper,top=2cm,bottom=2.2cm,left=1.8cm,right=1.8cm,
            columnsep=0.6cm]{geometry}
\usepackage{amsmath,amssymb,amsthm,mathtools}
\usepackage{bm}
\usepackage{booktabs}
\usepackage{array}
\usepackage{enumitem}
\usepackage{graphicx}
\usepackage{xcolor}
\usepackage{microtype}
\usepackage{algorithm}
\usepackage{algpseudocode}
\usepackage[round,sort&compress]{natbib}
\usepackage[colorlinks=true,linkcolor=blue!60!black,
            citecolor=blue!60!black,urlcolor=blue!60!black]{hyperref}
\usepackage{fancyhdr}

\pagestyle{fancy}
\fancyhf{}
\fancyhead[LE,RO]{\thepage}
\fancyhead[RE]{\small Categorical Search for Metabolomics Annotation}
\fancyhead[LO]{\small Sachikonye}
\renewcommand{\headrulewidth}{0.4pt}

% ---- Theorem environments -----------------------------------------------
\newtheoremstyle{thmstyle}{\topsep}{\topsep}{\itshape}{}{\bfseries}{.}{.5em}{}
\theoremstyle{thmstyle}
\newtheorem{theorem}{Theorem}[section]
\newtheorem{lemma}[theorem]{Lemma}
\newtheorem{proposition}[theorem]{Proposition}
\newtheorem{corollary}[theorem]{Corollary}
\theoremstyle{definition}
\newtheorem{definition}[theorem]{Definition}
\newtheorem{axiom}{Axiom}
\newtheorem{example}[theorem]{Example}
\theoremstyle{remark}
\newtheorem{remark}[theorem]{Remark}

% ---- Custom commands ----------------------------------------------------
\newcommand{\R}{\mathbb{R}}
\newcommand{\N}{\mathbb{N}}
\newcommand{\Z}{\mathbb{Z}}
\newcommand{\Zp}{\mathbb{Z}^{+}}
\newcommand{\kB}{k_{\mathrm{B}}}
\newcommand{\Parts}{\mathcal{P}}
\newcommand{\Sv}{\mathbf{S}}
\newcommand{\Sk}{S_{k}}
\newcommand{\St}{S_{t}}
\newcommand{\Se}{S_{e}}
\newcommand{\nP}{n_{\mathrm{P}}}
\newcommand{\calG}{\mathcal{G}}
\newcommand{\calV}{\mathcal{V}}
\newcommand{\calE}{\mathcal{E}}
\newcommand{\calT}{\mathcal{T}}
\newcommand{\calC}{\mathcal{C}}
\newcommand{\calL}{\mathcal{L}}
\newcommand{\calS}{\mathcal{S}}
\newcommand{\calR}{\mathcal{R}}
\newcommand{\calD}{\mathcal{D}}
\newcommand{\calH}{\mathcal{H}}
\newcommand{\flo}{\beta}
\DeclareMathOperator{\mean}{mean}
\DeclareMathOperator{\dist}{dist}
\DeclareMathOperator*{\argmin}{arg\,min}

% =========================================================
\title{\textbf{Categorical Search for Metabolomics Annotation:\\
Generative Phase-Space Addressing with Coherence-Grounded\\
Verification for Unknown Compound Identification}}

\author{
Kundai Farai Sachikonye\\
Technical University of Munich\\
Chair of Proteomics and Bioanalytics\\
\texttt{kundai.sachikonye@wzw.tum.de}
}
\date{\today}
% =========================================================

\begin{document}
\maketitle

% =========================================================
\begin{abstract}
\noindent
We present a complete search algorithm for the annotation of unknown
compounds from tandem mass spectrometry (MS/MS) data. The algorithm
operates in three stages. First, an observed spectrum is embedded into
three-dimensional S-entropy coordinate space $\Sv = (\Sk, \St, \Se)
\in [0,1]^3$ via a platform-independent transformation that encodes
mass-scale, structural complexity, and harmonic connectivity of the
fragment ion series. Second, the S-entropy triple is encoded as an
interleaved ternary address of depth~$k$, and a lazy trie over a
structural candidate database is traversed in $\mathcal{O}(k)$
operations to retrieve a candidate set of size
$|\calC_k| \ll N_{\mathrm{db}}$. Third, each candidate is verified
by three independent channels~--- spectral intensity prediction from
phase-lock edge density, fragmentation pathway admissibility from
partition-state graph search, and neutral loss consistency from
functional-group phase coherence~--- required to form a coherence
triangle before the annotation is accepted. The algorithm terminates
at closure: the state in which no further available verification
channel produces a materially different candidate, strictly stronger
than a confidence threshold. If the verification channels disagree,
the algorithm emits an explicit decline carrying the distinct candidate
classes, rather than silently selecting one.

We prove: (i)~the per-query cost is $\mathcal{O}(k + |\calC_k| \cdot
F)$ where $F$ is the per-candidate verification cost, independent of
database size $N_{\mathrm{db}}$ after a one-time $\mathcal{O}(N_{\mathrm{db}}
\cdot k)$ trie construction; (ii)~the ternary encoding preserves
chemical similarity, so that candidates sharing a longer address prefix
are metrically closer in S-entropy space; (iii)~the coherence triangle
requirement is both necessary and sufficient for robustness against
single-channel failure; and (iv)~every annotation carries an explicit
residual quantifying the irreducible gap between the observed spectrum
and any finite candidate set, derived from the non-completability of
the chemical medium.

\bigskip\noindent\textbf{Keywords:}
metabolomics annotation; database search; S-entropy coordinates;
ternary phase-space addressing; categorical fragmentation; coherence
triangle; partition-state graph; generative spectral database;
unknown compound identification; tandem mass spectrometry
\end{abstract}

\tableofcontents

% =========================================================
\section{Introduction}
\label{sec:intro}
% =========================================================

\subsection{The annotation problem}

Untargeted metabolomics by liquid chromatography coupled to tandem mass
spectrometry (LC-MS/MS) routinely detects thousands of molecular features
per sample, of which approximately 80\% remain
unannotated~\citep{fiehn2002metabolomics,kind2018identification}. The
annotation problem is: given a precursor mass-to-charge ratio $m/z$, a
charge state~$z$, a retention time, and a fragment ion spectrum
$\{(m/z_i, I_i)\}_{i=1}^{K}$, determine the molecular identity of the
precursor.

Two classes of approach dominate. \textbf{Spectral library matching}
scores an observed spectrum against a database of reference spectra
acquired on the same or similar instrument
types~\citep{stein2012mass,horai2010massbank}. The cost is
$\mathcal{O}(N_{\mathrm{db}} \times d)$ per query, where
$N_{\mathrm{db}}$ is the library size and $d$ is the spectral
dimensionality. The fundamental limitation is coverage: a compound
absent from the library cannot be identified, and compounds acquired on
a different instrument type may not match due to platform-dependent
intensity variations~\citep{kind2007fiehnlib}. \textbf{In-silico
fragmentation} predicts theoretical spectra for candidate structures
and scores against the observed spectrum~\citep{duhrkop2019sirius,
wei2019rapid}. This extends coverage to any structure whose
fragmentation can be simulated, but the per-candidate cost of quantum-
chemical or rule-based prediction is substantial, and the candidate
enumeration from molecular formula alone can exceed $10^6$ for masses
above 400\,Da.

Both approaches treat the spectrum as an empirical fingerprint. Neither
exploits the geometric structure of the fragmentation process: that
each ion occupies a specific location in a bounded oscillatory phase
space, and that fragmentations are transitions between locations in a
structured graph whose topology determines both which fragments appear
and at what relative intensities.

\subsection{Approach}

We present an algorithm that replaces the linear scan of spectral
library matching with a logarithmic-depth trie traversal, and replaces
the black-box scoring of in-silico methods with a three-channel
verification grounded in independently derived physical theories:

\begin{enumerate}[nosep,label=(\roman*)]
\item \textbf{S-entropy embedding and ternary addressing}
  (Section~\ref{sec:embedding}). Each spectrum is mapped to a
  platform-independent three-dimensional coordinate
  $\Sv = (\Sk, \St, \Se) \in [0,1]^3$ and encoded as a ternary string
  of depth~$k$. A candidate database pre-indexed as a ternary trie is
  traversed in $\mathcal{O}(k)$, returning a candidate set
  $\calC_k$ of compounds whose S-entropy coordinates share a common
  prefix with the query.

\item \textbf{Categorical intensity verification}
  (Section~\ref{sec:intensity}). For each candidate in $\calC_k$,
  the phase-lock edge density $|E_i|$ of every predicted fragment is
  computed from the candidate's molecular graph, yielding predicted
  intensities $I_i \propto \exp(-|E_i|/\langle E \rangle)$. The
  correlation between predicted and observed intensities is the first
  verification channel.

\item \textbf{Partition-state pathway verification}
  (Section~\ref{sec:pathway}). Each candidate's precursor is mapped to
  a partition state $(n, \ell, m, s)$ via the partition bijection, and
  the S-Entropy Bidirectional Dijkstra (SEBD-MS) algorithm is run to
  verify that all observed fragments lie within the forward
  reachability cone. The fraction of covered fragments is the second
  verification channel.

\item \textbf{Neutral loss phase-coherence verification}
  (Section~\ref{sec:neutralloss}). For each observed neutral loss, the
  candidate's molecular structure is checked for the corresponding
  functional group with residual phase coherence to the precursor. The
  fraction of correctly predicted neutral losses is the third
  verification channel.

\item \textbf{Coherence-grounded acceptance}
  (Section~\ref{sec:coherence}). The three channels must form a
  coherence triangle~--- each channel mutually supports the other
  two~--- before the annotation is accepted. Acceptance occurs at
  closure: the point where no further verification channel can produce
  a materially different candidate. If the channels disagree, the
  algorithm emits an explicit decline.
\end{enumerate}

\subsection{Relation to prior work}

The components of this algorithm have been developed independently.
The S-entropy coordinate system and ternary encoding were introduced
in~\citet{sachikonye2024generative} for generative spectral databases.
The partition-state graph and SEBD-MS algorithm were developed
in~\citet{sachikonye2024partition} for MS/MS fragment identification.
The categorical fragmentation theory and intensity model were
established in~\citet{sachikonye2024categorical}. The coherence
triangle and closure criterion were derived
in~\citet{sachikonye2024scp} as domain-agnostic search theory. The
platform-independent coordinate transformation was validated on
LC-MS/MS data in~\citet{sachikonye2024ucdavis}. This paper synthesises
these components into a single executable algorithm and proves its
correctness and complexity.

The novelty is the synthesis. No prior work combines ternary phase-space
addressing (sublinear candidate retrieval) with categorical intensity
prediction (first-principles scoring) and coherence-grounded acceptance
(principled stopping and decline). Each component addresses a different
failure mode of existing annotation pipelines:
\begin{itemize}[nosep]
\item Ternary addressing eliminates the $\mathcal{O}(N_{\mathrm{db}})$
  scan, making the algorithm database-size-independent at query time.
\item Categorical intensity prediction provides a mechanistic score
  without requiring training data, enabling annotation of novel
  chemical classes absent from any training set.
\item Coherence-grounded acceptance prevents silent misannotation: a
  result is never emitted unless three independent channels agree, and
  disagreement is reported explicitly rather than suppressed.
\end{itemize}


% =========================================================
\section{Mathematical Foundations}
\label{sec:foundations}
% =========================================================

\subsection{The Bounded Phase Space axiom}

\begin{axiom}[Bounded Phase Space]
\label{ax:bps}
Every persistent physical system occupies a bounded region
$\Omega \subset \R^{2d}$ of phase space with finite Liouville measure
$\mu(\Omega) < \infty$.
\end{axiom}

By the Poincar\'{e} recurrence theorem~\citep{poincare1890}, any
bounded Hamiltonian system without fixed points is oscillatory. For
integrable systems, the Arnold--Liouville theorem~\citep{arnold1989}
guarantees action-angle coordinates with quasi-periodic motion.

\subsection{Partition states and the bijection}

\begin{definition}[Partition State]
\label{def:partition-state}
A partition state of a bounded three-dimensional Hamiltonian system is
an equivalence class of phase-space points sharing closure label
$(n, \ell, m, s)$ where $n \in \{1,2,\ldots\}$ is the principal level,
$\ell \in \{0,\ldots,n-1\}$ the angular sublevel,
$m \in \{-\ell,\ldots,+\ell\}$ the azimuthal projection, and
$s \in \{-\tfrac{1}{2},+\tfrac{1}{2}\}$ the rotational parity. The
number of states at level~$n$ is $C(n) = 2n^2$, and the cumulative
count is $N_{\mathrm{state}}(n) = n(n+1)(2n+1)/3$.
\end{definition}

\begin{definition}[Partition Bijection $\Phi$]
\label{def:bijection}
The map $\Phi\colon \Zp \to \Parts$ sends each positive integer
$M$ to a unique partition state $(n,\ell,m,s)$ by: (i)~find $n$ with
$N_{\mathrm{state}}(n-1) < M \le N_{\mathrm{state}}(n)$; (ii)~set
$\delta = M - N_{\mathrm{state}}(n-1)$; (iii)~assign $(\ell,m,s)$ as
the $\delta$-th element of the lexicographic enumeration of level-$n$
states. $\Phi$ is a bijection~\citep{sachikonye2024partition}.
\end{definition}

An ion with mass-to-charge ratio $m/z$ maps to partition coordinates
via $n_{\mathrm{ion}} = \lfloor\sqrt{m/z}\rfloor + 1$.

\subsection{S-entropy coordinates}

\begin{definition}[Spectral S-Entropy Embedding]
\label{def:sentropy}
For a tandem mass spectrum with $K$ fragment peaks
$\{(m/z_i, I_i)\}_{i=1}^K$ and precursor $m/z_p$, the S-entropy
coordinates are:
\begin{align}
\Sk &= -\frac{1}{\ln K} \sum_{i=1}^{K} p_i \ln p_i,
  \quad p_i = \frac{I_i}{\sum_j I_j},
\label{eq:sk}\\[4pt]
\St &= \frac{\ln(m/z_{\max} / m/z_{\min})}
            {\ln(\omega_{\mathrm{ref,max}}/\omega_{\mathrm{ref,min}})},
\label{eq:st}\\[4pt]
\Se &= \frac{N_{\mathrm{harmonic}}}
            {\max\bigl(\binom{K}{2}, 1\bigr)},
\label{eq:se}
\end{align}
where $\Sk$ is the normalised Shannon entropy of the intensity
distribution (knowledge entropy), $\St$ is the normalised log-ratio
of the fragment $m/z$ range (temporal entropy), and $\Se$ is the
fraction of fragment pairs in harmonic proximity (evolution
entropy). Two fragments $m/z_a, m/z_b$ are in harmonic proximity if
there exist integers $p, q$ with $1 \le p, q \le 8$ such that
$|\max(m/z_a, m/z_b)/\min(m/z_a, m/z_b) - p/q| < \delta$ with
tolerance $\delta = 0.05$. Reference constants
$\omega_{\mathrm{ref,max}} = 4401\,\mathrm{cm}^{-1}$ and
$\omega_{\mathrm{ref,min}} = 67\,\mathrm{cm}^{-1}$ are the extremes
of known diatomic vibrational
frequencies~\citep{sachikonye2024generative}.
\end{definition}

\begin{proposition}[Well-definedness and Range]
\label{prop:sentropy-range}
For any spectrum with $K \ge 2$ peaks and positive intensities:
$\Sk \in [0,1]$, $\St \in [0,1]$, $\Se \in [0,1]$.
\end{proposition}

\begin{proof}
$\Sk$: Shannon entropy of a discrete distribution on $K$ outcomes is
maximised at $\ln K$ (uniform distribution) and minimised at $0$
(degenerate); normalisation by $\ln K$ maps to $[0,1]$. $\St$: the
numerator is $\ln(m/z_{\max}/m/z_{\min}) \ge 0$ and is bounded above
by the denominator by choice of reference constants spanning the full
physical range. $\Se$: a count of pairs divided by the total number
of pairs is in $[0,1]$ by definition. \qed
\end{proof}

\begin{remark}[Platform independence]
\label{rem:platform}
$\Sk$ depends on normalised intensities (ratios), not absolute
intensities; $\St$ depends on $m/z$ ratios; $\Se$ depends on $m/z$
ratios. All three are therefore invariant to absolute intensity
scaling and to systematic $m/z$ shifts within calibration tolerance.
This is the mathematical basis for the measured CV~$< 1.8\%$ across
Waters Q-TOF and Thermo Orbitrap
platforms~\citep{sachikonye2024categorical}.
\end{remark}


% =========================================================
\section{Ternary Phase-Space Addressing}
\label{sec:embedding}
% =========================================================

\subsection{Interleaved ternary encoding}

\begin{definition}[Interleaved Ternary Address]
\label{def:ternary}
For $\Sv = (\Sk, \St, \Se) \in [0,1]^3$ and depth $k$, the
interleaved ternary address $\alpha^{(k)}(\Sv) \in \{0,1,2\}^k$ is
computed by Algorithm~\ref{alg:encode}.
\end{definition}

\begin{algorithm}[h]
\caption{Interleaved Ternary Encoder}
\label{alg:encode}
\begin{algorithmic}[1]
\Require $\Sv = (\Sk,\St,\Se) \in [0,1]^3$; depth $k$
\Ensure Trit string $\alpha = (t_1,\ldots,t_k) \in \{0,1,2\}^k$
\State $r_0 \gets \Sk$;\; $r_1 \gets \St$;\; $r_2 \gets \Se$
\For{$j = 1$ \textbf{to} $k$}
  \State $d \gets (j-1) \bmod 3$
  \State $t_j \gets \min(\lfloor 3 r_d \rfloor, 2)$
  \State $r_d \gets 3 r_d - t_j$
\EndFor
\Return $(t_1, \ldots, t_k)$
\end{algorithmic}
\end{algorithm}

\begin{theorem}[Distance Preservation]
\label{thm:distance}
Two spectra with ternary addresses sharing a common prefix of length
$j$ satisfy:
\begin{equation}
\|\Sv_1 - \Sv_2\|_2 \le \sqrt{3} \cdot 3^{-\lfloor j/3 \rfloor}.
\label{eq:distance-bound}
\end{equation}
\end{theorem}

\begin{proof}
A common prefix of length $j$ means both addresses fall within the
same level-$j$ cell. After $\lfloor j/3 \rfloor$ complete refinement
cycles, each coordinate interval has width $3^{-\lfloor j/3 \rfloor}$.
The cell diameter is
$\sqrt{3} \cdot 3^{-\lfloor j/3 \rfloor}$. \qed
\end{proof}

\begin{corollary}[Chemical Similarity Preservation]
Compounds sharing a longer ternary prefix are closer in S-entropy
space. Chemical families (e.g., hydrogen halides, bent triatomics)
cluster by common prefix without any chemical knowledge being
encoded in the coordinate
system~\citep{sachikonye2024generative}.
\end{corollary}

\subsection{The lazy ternary trie}

\begin{definition}[Candidate Database Trie]
\label{def:trie}
Given a structural candidate database
$\calD = \{c_1, \ldots, c_{N_{\mathrm{db}}}\}$ with precomputed
S-entropy coordinates, the \emph{candidate trie} $\calT$ is a ternary
trie of depth $k$ where each compound $c_j$ is stored at the leaf
reached by following its ternary address $\alpha^{(k)}(\Sv_{c_j})$
from the root.
\end{definition}

\begin{theorem}[$\mathcal{O}(k)$ Candidate Retrieval]
\label{thm:retrieval}
Given a query spectrum with S-entropy triple $\Sv_q$, the candidate
set $\calC_k(\Sv_q) = \calT.\mathrm{prefix}(\alpha^{(k)}(\Sv_q))$
is retrieved in $\mathcal{O}(k)$ time, independent of
$N_{\mathrm{db}}$.
\end{theorem}

\begin{proof}
Each trie traversal step follows one of three child pointers (one per
trit value), performing $\mathcal{O}(1)$ work per step. After $k$
steps the leaf (or internal node for prefix queries) is reached. The
cost is $k$ pointer dereferences plus $\mathcal{O}(|\calC_k|)$ to
enumerate the candidates at the resolved cell. Since
$|\calC_k| \ll N_{\mathrm{db}}$ for $k \ge 12$ in
practice~\citep{sachikonye2024generative}, the total cost is
$\mathcal{O}(k + |\calC_k|)$. \qed
\end{proof}

\begin{remark}[Adaptive depth]
\label{rem:adaptive}
If the candidate set at depth $k$ is empty (no compound in the
database occupies that cell), the algorithm ascends to depth $k-1$,
broadening the search radius by a factor of~3 in the refined
dimension. This trades specificity for coverage and is particularly
relevant for unknown compounds not yet in any database: an empty cell
at depth $k$ with non-empty parent at depth $k-3$ predicts the
existence of a compound at the cell centre, as established by
the generative database framework~\citep{sachikonye2024generative}.
\end{remark}


% =========================================================
\section{Categorical Intensity Verification}
\label{sec:intensity}
% =========================================================

\subsection{The intensity model}

For each candidate $c \in \calC_k$, the molecular graph
$G_c = (V_c, E_c)$ is known (atoms as vertices, bonds as edges). The
fragmentation of $c$ under collision-induced dissociation produces a
set of predicted fragments, each inheriting a subgraph of $G_c$.

\begin{definition}[Phase-Lock Edge Density]
\label{def:edge-density}
For a fragment $f_i$ of candidate $c$ with molecular subgraph
$G_i = (V_i, E_i)$, the phase-lock edge density is
$|E_i| = |E(G_i)|$, the number of edges (bonds) in the fragment's
molecular graph. The reference edge count is
$\langle E \rangle = |E(G_c)|/N_{\mathrm{frag}}$, the mean edge count
per fragment.
\end{definition}

\begin{definition}[Predicted Fragment Intensity]
\label{def:predicted-intensity}
The predicted normalised intensity of fragment $f_i$ is:
\begin{equation}
\hat{I}_i = \frac{\alpha_i}{\sum_j \alpha_j},
\qquad
\alpha_i = \exp\!\left(-\frac{|E_i|}{\langle E \rangle}\right).
\label{eq:intensity}
\end{equation}
This is the topological entropy formulation: simple fragments (few
bonds, low $|E_i|$) have high termination probability and thus high
intensity; complex fragments (many bonds, high $|E_i|$) have low
intensity~\citep{sachikonye2024categorical}.
\end{definition}

\subsection{Spectral similarity score}

\begin{definition}[Intensity Verification Score]
\label{def:intensity-score}
The intensity verification score of candidate $c$ against observed
spectrum $\{(m/z_i, I_i^{\mathrm{obs}})\}$ is the Pearson correlation
between predicted and observed normalised intensities over matched
peaks:
\begin{equation}
\kappa_A(c) = \mathrm{corr}\!\bigl(
  \{\hat{I}_i\}_{i \in \calH(c)},\;
  \{I_i^{\mathrm{obs}}\}_{i \in \calH(c)}
\bigr),
\label{eq:kappa-a}
\end{equation}
where $\calH(c) = \{i : |\hat{m/z}_i - m/z_j^{\mathrm{obs}}| <
\varepsilon_{m/z} \text{ for some } j\}$ is the set of matched peaks
within mass tolerance $\varepsilon_{m/z}$.
\end{definition}

The correlation is computed without adjustable parameters beyond the
mass tolerance: $|E_i|$ and $\langle E \rangle$ are determined entirely
by the candidate's molecular graph. A correlation of $r = 0.87$ has
been validated across 2{,}847 MS/MS spectra on two
platforms~\citep{sachikonye2024categorical}.


% =========================================================
\section{Partition-State Pathway Verification}
\label{sec:pathway}
% =========================================================

\subsection{Forward reachability}

For each candidate $c$, the precursor ion maps to a partition state
$(n_p, \ell_p, m_p, s_p)$ via
$n_p = \lfloor\sqrt{m/z_p}\rfloor + 1$. Each observed fragment maps
similarly to $(n_f, \ell_f, m_f, s_f)$.

\begin{definition}[Forward Reachability Cone]
\label{def:forward-cone}
The forward reachability cone at depth $d_{\max}$ from precursor $n_p$
is:
\begin{equation}
\calR_f^{(d_{\max})} = \bigl\{v : \exists \text{ admissible path of
length } \le d_{\max} \text{ from } n_p \text{ to } v,\;
n_v \le n_p\bigr\}.
\end{equation}
An admissible transition $(u, v)$ requires mass reduction
($n_v \le n_u$), energy conservation, and charge conservation. The
size is bounded: $|\calR_f^{(d_{\max})}| \le 2n_p d_{\max} +
d_{\max}^2$~\citep{sachikonye2024partition}.
\end{definition}

\begin{definition}[Pathway Verification Score]
\label{def:pathway-score}
The pathway verification score is the fraction of observed fragments
falling within the forward reachability cone:
\begin{equation}
\kappa_B(c) = \frac{|\{i : n_{f_i} \in \calR_f^{(d_{\max})}(c)\}|}{K}.
\label{eq:kappa-b}
\end{equation}
\end{definition}

Validation on NIST HCD spectra yields
$\kappa_B = 0.947$~\citep{sachikonye2024partition}, with the remaining
5.3\% attributable to isotope peaks and adducts falling outside the
mass-reduction constraint.

\subsection{Virtual substate transition states}

The SEBD-MS algorithm searches backward from observed fragments through
virtual substate decompositions. A virtual predecessor
$\Sv^* = 2\Sv_f - \Sv_2$ for $\Sv_2 \in [0,1]^3$ may lie outside
$[0,1]^3$ (off-shell). Off-shell components correspond to
fragmentation transition states: intermediate configurations that
mediate bond cleavage but cannot be isolated as stable
ions~\citep{sachikonye2024partition}. The off-shell fraction
(8.3\% in validation) provides additional structural information:
candidates whose predicted off-shell fraction matches the observed
fraction receive a bonus to $\kappa_B$.


% =========================================================
\section{Neutral Loss Phase-Coherence Verification}
\label{sec:neutralloss}
% =========================================================

\subsection{Functional group phase coherence}

Neutral losses ($-18$\,Da for H$_2$O, $-44$\,Da for CO$_2$,
$-17$\,Da for NH$_3$) are not independent stochastic events but
phase-correlated processes arising from residual correlations between
fragments and precursor functional
groups~\citep{sachikonye2024categorical}.

\begin{definition}[Phase Coherence]
\label{def:phase-coherence}
For a candidate $c$ with functional group $g$ (e.g., hydroxyl,
carboxyl, amine), the phase coherence between the precursor and the
fragment retaining the functional group's site is:
\begin{equation}
C_{\phi}(g) = \bigl|\bigl\langle
  e^{i(\phi_g - \phi_{\mathrm{prec}})}
\bigr\rangle\bigr|,
\label{eq:phase-coherence}
\end{equation}
where the phase $\phi_g$ is determined by the fragment's position in
S-entropy space and $\phi_{\mathrm{prec}}$ by the precursor's
position. In practice, $C_{\phi}(g)$ is computed from the S-entropy
distance between the fragment and the precursor projected onto the
$\Se$ axis (which encodes harmonic connectivity).
\end{definition}

\begin{definition}[Neutral Loss Verification Score]
\label{def:nl-score}
The neutral loss verification score is the fraction of observed neutral
losses correctly predicted by the candidate's functional group
inventory with phase coherence exceeding threshold $\theta_{\phi}$:
\begin{equation}
\kappa_C(c) = \frac{|\{l \in \calL_{\mathrm{obs}} :
  \exists\, g \in \mathrm{FG}(c),\;
  \mathrm{mass}(g) = l,\;
  C_{\phi}(g) > \theta_{\phi}\}|}
  {|\calL_{\mathrm{obs}}|},
\label{eq:kappa-c}
\end{equation}
where $\calL_{\mathrm{obs}}$ is the set of observed neutral losses,
$\mathrm{FG}(c)$ is the functional group inventory of candidate $c$,
and $\theta_{\phi} = 0.3$ is the coherence
threshold~\citep{sachikonye2024categorical}.
\end{definition}

Validation yields $P(\text{loss} \mid C_{\phi} > 0.3) = 0.943 \pm
0.032$ versus $P(\text{loss} \mid C_{\phi} \le 0.3) = 0.247 \pm
0.071$, a 3.8-fold enhancement confirming the phase-memory
mechanism~\citep{sachikonye2024categorical}.


% =========================================================
\section{Coherence-Grounded Acceptance}
\label{sec:coherence}
% =========================================================

\subsection{The coherence triangle}

The three verification scores
$(\kappa_A, \kappa_B, \kappa_C)$~--- intensity, pathway, neutral
loss~--- are computed from independent physical models (topological
entropy, partition-state graph, phase coherence respectively). We now
require that they form a coherence triangle before accepting an
annotation.

\begin{definition}[Support Relation]
\label{def:support}
Verification channel $j$ \emph{supports} channel $i$ at threshold
$\theta$ if the presence of channel $j$'s score does not reduce, and
to degree $\ge \theta$ reinforces, channel $i$'s discrimination between
the top-ranked candidate and the second-ranked candidate. The support
graph $G_S$ has an edge $j \to i$ whenever $j$ supports $i$.
\end{definition}

\begin{theorem}[Coherence Triangle Necessity]
\label{thm:triangle}
If an annotation is to be robust against the failure of any single
verification channel, the support graph $G_S$ must contain a directed
cycle of length $\ge 3$. No linear chain and no pair of mutually
supporting channels suffices~\citep{sachikonye2024scp}.
\end{theorem}

\begin{proof}
A linear support chain has an unsupported terminal channel; its removal
collapses the justification. A 2-cycle is mutual definition between
two channels with no independent third check. Only a cycle of length
$\ge 3$ provides each channel with at least two independent supporters,
so that removal of any one leaves the remaining two still mutually
supporting. With three channels at mutual support threshold
$\theta > \tfrac{1}{2}$, the two survivors jointly exceed any single
dissenting signal. \qed
\end{proof}

\begin{definition}[Composite Verification Power]
\label{def:composite}
The composite verification power of the three channels is:
\begin{equation}
\kappa_{\mathrm{comp}} = 1 - (1 - \kappa_A)(1 - \kappa_B)(1 - \kappa_C).
\label{eq:composite}
\end{equation}
This is the multiplicative composition law: the fraction of the
initial above-floor gap closed by the combined verification, with
each channel closing its own fraction
independently~\citep{sachikonye2024scp}.
\end{definition}

\subsection{The resolution floor}

\begin{definition}[Resolution Floor]
\label{def:floor}
The resolution floor $\flo > 0$ is the irreducible residual between
any finite candidate set and the true identity of an unknown compound.
It follows from the non-completability of the chemical medium: for any
finite database, there exist compounds not in the database, so no
finite search exhausts the
medium~\citep{sachikonye2024scp,sachikonye2024generative}.
\end{definition}

\begin{theorem}[Every Annotation Carries a Residual]
\label{thm:residual}
For every accepted annotation with composite verification power
$\kappa_{\mathrm{comp}}$, the residual above-floor gap is:
\begin{equation}
r = (1 - \kappa_{\mathrm{comp}})(a_0 - \flo),
\label{eq:residual}
\end{equation}
where $a_0$ is the initial alignment gap. This residual is strictly
positive for all finite verification efforts: no annotation achieves
$r = 0$.
\end{theorem}

\begin{proof}
By the multiplicative law
(Equation~\ref{eq:composite}), each verification channel closes a
fraction of the remaining gap. The product
$(1-\kappa_A)(1-\kappa_B)(1-\kappa_C) > 0$ whenever each
$\kappa_i < 1$. Since no finite verification can close the full gap
(this would require exhaustive comparison against the entire chemical
medium, which is non-completable), $r > 0$. \qed
\end{proof}

\subsection{Closure criterion}

\begin{definition}[Closure]
\label{def:closure}
The annotation search for a given query spectrum is \emph{closed} if,
for every verification channel not yet applied and every candidate not
yet tested, extending the search cannot add a candidate to the
accepted set that is materially different from the current best
candidate. Formally, closure holds when:
\begin{equation}
\forall\, c' \in \calC_k \setminus \{c^*\}:\quad
\kappa_{\mathrm{comp}}(c') < \kappa_{\mathrm{comp}}(c^*) - \flo,
\label{eq:closure}
\end{equation}
where $c^*$ is the current best candidate and $\flo$ is the resolution
floor.
\end{definition}

\begin{theorem}[Closure is Strictly Stronger Than a Confidence Threshold]
\label{thm:closure-stronger}
There exist spectra and databases for which
$\kappa_{\mathrm{comp}}(c^*) > \theta$ for any fixed threshold
$\theta$, yet the search is not closed: a further untested candidate
$c'$ may achieve $\kappa_{\mathrm{comp}}(c') > \kappa_{\mathrm{comp}}(c^*)$~\citep{sachikonye2024scp}.
\end{theorem}

\begin{definition}[Decline]
\label{def:decline}
If closure is reached but two or more candidates have composite
verification power within $\flo$ of each other, the algorithm emits a
\emph{decline}: an explicit statement that the evidence is insufficient
to distinguish the candidates, together with the distinct candidate
classes found. A decline is a first-class output, not an error.
\end{definition}


% =========================================================
\section{The Complete Algorithm}
\label{sec:algorithm}
% =========================================================

\begin{algorithm}[t]
\caption{Categorical Search for Metabolomics Annotation}
\label{alg:main}
\begin{algorithmic}[1]
\Require Observed spectrum $\{(m/z_i, I_i)\}_{i=1}^K$; precursor
  $m/z_p$; charge $z$; candidate trie $\calT$; depth $k$;
  floor $\flo$; coherence threshold $\theta$; mass tolerance
  $\varepsilon_{m/z}$; phase coherence threshold $\theta_{\phi}$
\Ensure Annotation $c^*$ with verification scores, or decline with
  candidate classes
\Statex
\State \textbf{// Stage 1: S-Entropy Embedding}
\State Compute $\Sk, \St, \Se$ via
  Equations~\eqref{eq:sk}--\eqref{eq:se}
\State $\Sv_q \gets (\Sk, \St, \Se)$
\Statex
\State \textbf{// Stage 2: Ternary Candidate Retrieval}
\State $\alpha \gets \textsc{TernaryEncode}(\Sv_q, k)$
  \Comment{Algorithm~\ref{alg:encode}, $\mathcal{O}(k)$}
\State $\calC_k \gets \calT.\mathrm{prefix}(\alpha)$
  \Comment{$\mathcal{O}(k)$}
\While{$\calC_k = \emptyset$ and $k > 3$}
  \Comment{Adaptive depth}
  \State $k \gets k - 3$;\quad
    $\alpha \gets \alpha[1..k]$;\quad
    $\calC_k \gets \calT.\mathrm{prefix}(\alpha)$
\EndWhile
\If{$\calC_k = \emptyset$}
  \State \Return \textsc{Decline}(``no candidates at any depth'')
\EndIf
\Statex
\State \textbf{// Stage 3: Three-Channel Verification}
\For{each candidate $c \in \calC_k$}
  \State Compute $\kappa_A(c)$ via Equation~\eqref{eq:kappa-a}
    \Comment{Intensity}
  \State Compute $\kappa_B(c)$ via Equation~\eqref{eq:kappa-b}
    \Comment{Pathway}
  \State Compute $\kappa_C(c)$ via Equation~\eqref{eq:kappa-c}
    \Comment{Neutral loss}
  \State $\kappa_{\mathrm{comp}}(c) \gets
    1-(1-\kappa_A)(1-\kappa_B)(1-\kappa_C)$
\EndFor
\Statex
\State \textbf{// Stage 4: Coherence Check}
\State $c^* \gets \argmin_{c \in \calC_k}
  (1 - \kappa_{\mathrm{comp}}(c))$
\State Build support graph $G_S$ over $(\kappa_A, \kappa_B, \kappa_C)$
  for $c^*$
\If{$G_S$ has no cycle of length $\ge 3$}
  \State Emit \textsc{CoherenceWarning}
\EndIf
\Statex
\State \textbf{// Stage 5: Closure or Decline}
\State Sort $\calC_k$ by $\kappa_{\mathrm{comp}}$ descending
\If{$\kappa_{\mathrm{comp}}(c^*) -
    \kappa_{\mathrm{comp}}(c^*_2) > \flo$}
  \Comment{Closure}
  \State $r \gets (1-\kappa_{\mathrm{comp}}(c^*))(1 - \flo)$
  \State \Return $(c^*, \kappa_A, \kappa_B, \kappa_C, r)$
\Else
  \Comment{Contested}
  \State \Return \textsc{Decline}($\{c :
    \kappa_{\mathrm{comp}}(c) \ge
    \kappa_{\mathrm{comp}}(c^*) - \flo\}$)
\EndIf
\end{algorithmic}
\end{algorithm}


% =========================================================
\section{Complexity Analysis}
\label{sec:complexity}
% =========================================================

\begin{theorem}[Total Query Complexity]
\label{thm:complexity}
Algorithm~\ref{alg:main} has total per-query cost:
\begin{equation}
T_{\mathrm{query}} = \underbrace{\mathcal{O}(K)}_{\text{S-entropy}} +
  \underbrace{\mathcal{O}(k)}_{\text{trie}} +
  \underbrace{\mathcal{O}(|\calC_k| \cdot F)}_{\text{verification}},
\label{eq:complexity}
\end{equation}
where $K$ is the number of fragment peaks, $k$ is the trie depth, and
$F$ is the per-candidate verification cost. The verification cost is:
\begin{equation}
F = \underbrace{\mathcal{O}(K)}_{\kappa_A} +
  \underbrace{\mathcal{O}(n_p d_{\max} \log(n_p d_{\max}))}_{\kappa_B} +
  \underbrace{\mathcal{O}(|\calL_{\mathrm{obs}}| \cdot
    |\mathrm{FG}(c)|)}_{\kappa_C}.
\label{eq:F}
\end{equation}
\end{theorem}

\begin{proof}
Stage~1 computes three scalar functions of $K$ peaks: $\mathcal{O}(K)$.
Stage~2 is trie traversal: $\mathcal{O}(k)$
(Theorem~\ref{thm:retrieval}). Stage~3 applies three verification
functions to each of $|\calC_k|$ candidates. $\kappa_A$ requires
matching $K$ observed peaks against predicted peaks and computing
a correlation: $\mathcal{O}(K)$. $\kappa_B$ runs the forward
reachability computation of SEBD-MS: $\mathcal{O}(n_p d_{\max}
\log(n_p d_{\max}))$~\citep{sachikonye2024partition}. $\kappa_C$
checks each observed neutral loss against each functional group:
$\mathcal{O}(|\calL_{\mathrm{obs}}| \cdot |\mathrm{FG}(c)|)$. Stages~4
and~5 are $\mathcal{O}(|\calC_k| \log |\calC_k|)$ (sorting). The
dominant term is Stage~3. \qed
\end{proof}

\begin{corollary}[Database-Size Independence]
\label{cor:db-independence}
The per-query cost is independent of $N_{\mathrm{db}}$. The database
size affects only the one-time trie construction cost
$\mathcal{O}(N_{\mathrm{db}} \cdot k)$.
\end{corollary}

\begin{table}[h]
\centering
\caption{Complexity comparison with existing approaches.}
\label{tab:complexity}
\small
\begin{tabular}{@{}lcc@{}}
\toprule
\textbf{Method} & \textbf{Per-Query Cost} & \textbf{Depends on $N_{\mathrm{db}}$?} \\
\midrule
Spectral library search & $\mathcal{O}(N_{\mathrm{db}} \cdot d)$ & Yes \\
Fingerprint search & $\mathcal{O}(N_{\mathrm{db}} \cdot d_{\mathrm{fp}})$ & Yes \\
In-silico fragmentation & $\mathcal{O}(|\mathrm{cand}| \cdot T_{\mathrm{sim}})$ & Yes$^\dagger$ \\
Categorical search (ours) & $\mathcal{O}(k + |\calC_k| \cdot F)$ & No \\
\bottomrule
\end{tabular}

\smallskip
{\footnotesize $^\dagger$Candidate enumeration from molecular formula
scales with database coverage.}
\end{table}

\begin{corollary}[Speedup at Scale]
\label{cor:speedup}
For $N_{\mathrm{db}} = 10^8$ (PubChem scale), $k = 18$,
$|\calC_k| = 10$, and $F = 10^3$ operations per candidate: the
categorical search cost is $\approx 10^4$ operations per query versus
$\approx 10^{11}$ for exhaustive spectral library search, a speedup
of $\sim 10^7$.
\end{corollary}


% =========================================================
\section{CASMI-2022 Application}
\label{sec:casmi}
% =========================================================

The CASMI (Critical Assessment of Small Molecule Identification) 2022
challenge~\citep{fiehn2022casmi} provides 500 unknown compounds
acquired by LC-MS/MS, deliberately excluded from public spectral
libraries, with evaluation on four categories: (1)~adduct annotation,
(2)~elemental formula, (3)~compound structure class, and
(4)~2D chemical structure.

\subsection{Mapping to Algorithm Outputs}

\textbf{Adduct annotation} (Category~1). The partition bijection maps
$m/z_p$ to $n_p = \lfloor\sqrt{m/z_p}\rfloor + 1$. The charge state
$z$ is determined by the partition malformation energy:
$\mathcal{M}_{\mathrm{ion}}(Z, z) = \kB T_{\mathrm{cat}}
\ln(Z!/(Z-z)!)$~\citep{sachikonye2024partition}. The adduct type
($[\mathrm{M}+\mathrm{H}]^+$, $[\mathrm{M}+\mathrm{Na}]^+$, etc.)
is identified by comparing the observed $n_p$ against the predicted
$n_p$ for each adduct hypothesis.

\textbf{Elemental formula} (Category~2). The candidate's molecular
formula is a direct output of the trie lookup: each candidate in
$\calC_k$ carries its molecular formula.

\textbf{Compound class} (Category~3). The ternary prefix at coarse
depth ($k = 3$) determines the compound class. Compounds sharing a
3-trit prefix co-locate in S-entropy space without chemical knowledge
being encoded: the cell $[202]$ captures molecules with high $\Sk$,
low $\St$, high $\Se$~\citep{sachikonye2024generative}.

\textbf{2D structure} (Category~4). The full-depth ($k = 12$--$18$)
trie resolution, combined with the three-channel verification, selects
the best structural candidate. The output is the candidate's InChIKey
or SMILES.

\subsection{Handling unknown compounds}

For compounds not in any structural database, the algorithm's adaptive
depth mechanism (Remark~\ref{rem:adaptive}) broadens the search. If no
candidate is found at any depth, the S-entropy coordinates themselves
constrain the possible molecular properties:
\begin{itemize}[nosep]
\item $\Sk$ constrains the intensity distribution and hence the
  fragmentation complexity.
\item $\St$ constrains the $m/z$ range and hence the mass.
\item $\Se$ constrains the harmonic connectivity and hence the
  structural family.
\end{itemize}
Combined with the accurate mass and isotope pattern, these constraints
narrow the candidate space even without a database hit. The algorithm
reports a decline with the constraints as payload rather than forcing
an incorrect annotation.


% =========================================================
\section{Discussion}
\label{sec:discussion}
% =========================================================

\subsection{The circularity and its resolution}

The intensity model (Section~\ref{sec:intensity}) requires the
candidate's molecular graph to compute $|E_i|$. For an unknown
compound, the molecular graph is not known a priori. This is resolved
by the two-stage architecture: Stage~2 (ternary retrieval) reduces the
candidate set to $|\calC_k|$ compounds whose molecular graphs \emph{are}
known (they are entries in a structural database). Stage~3 then
computes $|E_i|$ from each candidate's known graph. The intensity
model is thus applied as a \emph{verification} step, not a generation
step: it verifies whether a known candidate's predicted spectrum
matches the observed spectrum, rather than generating a molecular graph
from an intensity pattern.

For compounds genuinely absent from any structural database (the
hardest CASMI cases), the intensity model is unavailable and the
algorithm falls back to $\kappa_B$ (pathway) and $\kappa_C$ (neutral
loss) verification only. The coherence triangle is then broken (only
two channels), and the algorithm emits a coherence warning rather than
a confident annotation. This is the honest outcome: with two channels,
the annotation is not robustly grounded against single-channel failure
(Theorem~\ref{thm:triangle}).

\subsection{Scope and limitations}

\textbf{Database dependence.} The algorithm's power depends on the
structural database indexing the trie. A trie built over PubChem
($\sim 10^8$ compounds) covers a vast chemical space; a trie over a
specialised metabolite database covers less space but with higher
relevance. The choice of database is a user decision, not a property of
the algorithm.

\textbf{S-entropy degeneracy.} At low trie depth, multiple chemically
distinct compounds may share a ternary cell (degeneracy). The three-channel
verification resolves this degeneracy: compounds with similar S-entropy
coordinates but different molecular graphs will have different
$\kappa_A$ scores. Full resolution is achieved at depth $k = 11$ for
39 compounds~\citep{sachikonye2024generative}; the required depth for
$10^8$ compounds is an open empirical question.

\textbf{Fragmentation coverage.} The intensity model assumes
collision-induced dissociation of singly charged small molecules
($m < 1500$\,Da). Extension to multi-charge states, electron-transfer
dissociation, and large biomolecules requires modifications to the
phase-lock network accounting~\citep{sachikonye2024categorical}.


% =========================================================
\section{Conclusion}
\label{sec:conclusion}
% =========================================================

We have presented a complete algorithm for metabolomics annotation that
synthesises four independently derived theoretical frameworks:
S-entropy coordinates for platform-independent spectral embedding,
ternary phase-space addressing for sublinear candidate retrieval,
categorical fragmentation theory for first-principles intensity
prediction, and coherence-grounded verification for principled
acceptance and decline.

The algorithm's five structural properties distinguish it from existing
annotation pipelines:

\begin{enumerate}[nosep,label=(\roman*)]
\item \textbf{Database-size independence}: per-query cost is
  $\mathcal{O}(k + |\calC_k| \cdot F)$, independent of
  $N_{\mathrm{db}}$.
\item \textbf{Training-data independence}: the intensity model has no
  adjustable parameters beyond the candidate's molecular graph.
\item \textbf{Platform independence}: S-entropy coordinates achieve
  CV~$< 1.8\%$ across instrument types.
\item \textbf{Principled acceptance}: three independent verification
  channels must form a coherence triangle before annotation is accepted.
\item \textbf{Honest decline}: disagreement between channels is
  reported explicitly as a first-class output.
\end{enumerate}

Every component reduces to a computation on finite weighted graphs or
finite sets: S-entropy coordinates are scalar functions of peak lists,
ternary addresses are string encodings, the partition-state graph is
a finite directed graph, and the coherence triangle is a support graph
on three nodes. The algorithm is therefore fully specified and directly
implementable.


% =========================================================
% REFERENCES
% =========================================================

\begin{thebibliography}{99}

\bibitem[Arnold(1989)]{arnold1989}
V.~I. Arnold.
\newblock \emph{Mathematical Methods of Classical Mechanics}.
\newblock Springer, 2nd edition, 1989.

\bibitem[D\"{u}hrkop et~al.(2019)]{duhrkop2019sirius}
K.~D\"{u}hrkop, M.~Fleischauer, M.~Ludwig, A.~A. Aksenov, A.~V. Melnik,
  M.~Meusel, P.~C. Dorrestein, J.~Rousu, and S.~B\"{o}cker.
\newblock {SIRIUS} 4: a rapid tool for turning tandem mass spectra into
  metabolite structure information.
\newblock \emph{Nature Methods}, 16:299--302, 2019.

\bibitem[Fiehn(2002)]{fiehn2002metabolomics}
O.~Fiehn.
\newblock Metabolomics---the link between genotypes and phenotypes.
\newblock \emph{Plant Molecular Biology}, 48:155--171, 2002.

\bibitem[Fiehn et~al.(2022)]{fiehn2022casmi}
O.~Fiehn et~al.
\newblock {CASMI} 2022: Critical Assessment of Small Molecule Identification.
\newblock UC Davis West Coast Metabolomics Center, 2022.
\newblock \url{https://fiehnlab.ucdavis.edu/casmi}

\bibitem[Horai et~al.(2010)]{horai2010massbank}
H.~Horai, M.~Arita, S.~Kanaya, Y.~Nihei, T.~Ikeda, K.~Suwa, Y.~Ojima,
  K.~Tanaka, S.~Tanaka, K.~Aoshima, et~al.
\newblock {MassBank}: a public repository for sharing mass spectral data for
  life sciences.
\newblock \emph{Journal of Mass Spectrometry}, 45:703--714, 2010.

\bibitem[Kind and Fiehn(2007)]{kind2007fiehnlib}
T.~Kind and O.~Fiehn.
\newblock Seven golden rules for heuristic filtering of molecular formulas
  obtained by accurate mass spectrometry.
\newblock \emph{BMC Bioinformatics}, 8:105, 2007.

\bibitem[Kind et~al.(2018)]{kind2018identification}
T.~Kind, H.~Tsugawa, T.~Cajka, Y.~Ma, Z.~Lai, S.~S. Mehta, G.~Wohlgemuth,
  D.~K. Barupal, M.~R. Showalter, M.~Arita, and O.~Fiehn.
\newblock Identification of small molecules using accurate mass {MS/MS}
  search.
\newblock \emph{Mass Spectrometry Reviews}, 37:513--532, 2018.

\bibitem[Poincar\'{e}(1890)]{poincare1890}
H.~Poincar\'{e}.
\newblock Sur le probl\`{e}me des trois corps et les \'{e}quations de la
  dynamique.
\newblock \emph{Acta Mathematica}, 13:1--270, 1890.

\bibitem[Sachikonye(2024a)]{sachikonye2024generative}
K.~F. Sachikonye.
\newblock Context-based spectral database: Generative phase space addressing
  for mass spectrometric identification from first principles.
\newblock Technical report, AIMe Registry for Artificial Intelligence, 2024.

\bibitem[Sachikonye(2024b)]{sachikonye2024partition}
K.~F. Sachikonye.
\newblock Partition-state graph search for tandem mass spectrometry:
  Bidirectional {D}ijkstra with virtual substate transition states.
\newblock Technical report, Technical University of Munich, 2024.

\bibitem[Sachikonye(2024c)]{sachikonye2024categorical}
K.~F. Sachikonye.
\newblock Categorical fragmentation networks in tandem mass spectrometry:
  Phase-lock topology and entropy-intensity relations.
\newblock Technical report, Technical University of Munich, 2024.

\bibitem[Sachikonye(2024d)]{sachikonye2024scp}
K.~F. Sachikonye.
\newblock Semantic causal propagation: An individuation-theoretic calculus of
  boundary-free search.
\newblock Technical report, Technical University of Munich, 2024.

\bibitem[Sachikonye(2024e)]{sachikonye2024ucdavis}
K.~F. Sachikonye.
\newblock Coordinate transformation for platform-independent metabolomics
  analysis: Application to the {UC Davis} dataset.
\newblock Technical report, Lavoisier Project, 2024.

\bibitem[Sachikonye(2024f)]{sachikonye2024gpu}
K.~F. Sachikonye.
\newblock Fragment shader as observation apparatus: {O}(1)-memory universal
  computation through the rendering--measurement identity.
\newblock Technical report, AIMe Registry for Artificial Intelligence, 2024.

\bibitem[Stein(2012)]{stein2012mass}
S.~E. Stein.
\newblock Mass spectral reference libraries: An ever-expanding resource for
  chemical identification.
\newblock \emph{Analytical Chemistry}, 84:7274--7282, 2012.

\bibitem[Wei et~al.(2019)]{wei2019rapid}
J.~N. Wei, D.~Duvenaud, and A.~Aspuru-Guzik.
\newblock Neural networks for the prediction of organic chemistry reactions.
\newblock \emph{ACS Central Science}, 2:725--732, 2019.

\end{thebibliography}

\end{document}